\scanpage{71}
\mspage{64}
% archive sheet number: 70
Ensuite, on voit séparément pour $h_1,h_2$ qu'on peut les trouver en \unsure{satisfaisant|astreignant} chacun aux deux contraintes de commutativité (1 triangle vertical et un carré penché). Pour $h_1$, p.ex., \unsure{l'existence} se voit sur le diagramme carré auxiliaire (traits pleins)
\diagram{
\begin{tikzcd}[ampersand replacement=\&,column sep=large,row sep=large]
A_1\arrow[r,"\mathrm{TC}","\deleted{\unsure{i}}"']\arrow[dd,dash] \& B_1\arrow[dd]\arrow[ddl,dashed,"h_1"']\arrow[dr,dashed,"\alpha_1^B"] \& \\
\&\&B_0\arrow[dl,dashed,"h_0"]\\
X_1\arrow[r,"F","\alpha_1^X"']\&X_0\&
\end{tikzcd}
}
\uline{NB} C'est la dém. de la prop. 3 p. 50

Le corollaire (partie \unsure{``il faut''}) étant vu, le \deleted{corollaire} lemme en résulte, car on sait que si on a un hom. $f$ de carrés cartésiens $X_\bullet$, $Y_\bullet$
\diagram{
\begin{tikzcd}[ampersand replacement=\&,column sep=large,row sep=large]
\&X_1\arrow[dl,"p_1^X"']\arrow[dd,"f_1"']\&\&X_3\arrow[ll]\arrow[dl]\arrow[dd,"f_3"]\\
X_0\arrow[dd,"f_0"']\&\&X_2\arrow[ll,"p_2^X"]\arrow[dd,"f_2"]\&\\
\&Y_1\arrow[dl,"p_1^Y"']\&\&Y_3\arrow[ll]\arrow[dl]\\
Y_0\&\&Y_2\arrow[ll,"p_2^Y"]\&
\end{tikzcd}
}
avec $p_{\unsure{1|2}}^X,p_2^X,p_1^Y,p_2^Y$ \unsure{fibrants|fibrantes} (il suffit $=$ $p_1^X$ et $p_1^Y$ \unsure{fibrants|fibrantes}), et $f_0,f_1,f_2\in W$ alors $f_3\in W$.

(\unsure{Ça|Ca} a été vu dans la théorie des produits fibrés $W$-homotopiques, et \deleted{\unsure{la}} vérification
